\section{KNN 把度量直接变成归纳偏置}
\label{sec:13-Machine-Learning/03-Trees-Ensembles-and-Neighbors/03}

你手里有一堆二手房交易记录。每套房记了面积、楼层、到地铁的步行分钟数，还有成交价。现在中介推来一套新房源，特征都有了，就缺价格。你打算怎么猜？

最直接的想法：从历史记录里找出和这套房最像的那几套，看看它们卖了多少钱，取个平均。不需要先训练一个全局公式，不需要假设价格和面积之间是直线还是曲线——决定预测结果的，就是历史数据里谁离查询点最近。

这就是 K-nearest neighbors。它几乎没有显式训练：给定查询 \(x\)，找到训练集中距离最近的 \(k\) 个点，回归取响应平均，分类投票。规则简单到一句话就能说完，但它把关于相似性的全部判断压进了两个地方：你用什么距离，以及特征长什么样。

\subsection{邻域内平均是在估计局部条件均值}

回归 KNN 的式子看着很朴素：

\[
\widehat f_k(x)
=\frac1k\sum_{i\in\mathcal N_k(x)}y_i,
\]

其中 \(\mathcal N_k(x)\) 是 \(x\) 的 \(k\) 个最近邻。写的是一串数据点的平均，估计的却是 \(\E[Y\mid X\approx x]\)——在查询点附近，响应变量的条件均值。

要做到这件事，需要两个条件同时成立。第一，邻域要足够小，落在里面的样本近似来自 \(X=x\) 附近，这样用它们的 \(y\) 去逼近 \(x\) 处的条件均值才有道理。第二，邻域里的样本要足够多，平均操作才能压低噪声。

麻烦的是这两个要求互相打架。想让邻域变小，就得减小 \(k\)；想让样本变多，就得增大 \(k\)。一致性要求两边妥协：\(k\to\infty\) 以汇聚更多样本，同时 \(k/n\to0\) 以保证邻域半径随着数据增多而收缩。

\(k\) 很小的时候，邻域只包住查询点附近极窄的一小片，局部细节保留得好，但几个点的小样本均值方差很大——换一组训练数据，最近邻可能就换了一批，预测跟着跳。\(k\) 很大时，平均涉及大片区域，输出很稳定，但远方的点也被拉进来了，把本不该混在一起的区域一锅炖。选 \(k\) 不是在挑投票人数，是在选平滑尺度。

\subsection{遭遇第一堵墙：谁算邻居由谁说了算}

Euclidean 距离公式写着

\[
d(x,x')^2
=\sum_{j=1}^{p}(x_j-x_j')^2.
\]

每个坐标平等地贡献平方差。但现实数据里坐标从来不平等。面积用平方米记，楼层就是整数，到地铁的距离用米记，三个坐标的量纲和取值范围天差地别。把到地铁的距离从米换成毫米，这一项在平方和里的贡献瞬间放大一百万倍；最近邻关系可能被距离的一个坐标完全绑架。

标准化是第一个补救：各坐标减均值除标准差，让每个特征的数值范围拉到相近。但这只是把问题从量纲转移到了方差——方差大的特征是否更重要的判断，交给了数据收集者当初选单位时的随意决定，而不是交给你要解决的那个问题。

Mahalanobis 距离往前走了一步：除以协方差矩阵，把各坐标的相关结构也考虑进去。学习型距离度量则走得更远，直接让数据或任务目标来决定什么叫相似——相当于把距离函数也当成了模型参数。KNN 的``训练''如果存在，就是学习一个好的距离。

\subsection{无关特征在暗处堆叠噪声}

更阴险的问题是无关特征。假设二十个坐标里只有三个和房价有关。剩下十七个坐标充满噪声：不论查询点取在哪，噪声坐标上的平方差照加不误。单个噪声坐标贡献不大，但十七个无意义的平方差累在一起，可能把三个有用特征发出的信号完全淹没。

你是在用全部坐标算距离，但只有部分坐标定义了真正的相似性。特征选择和度量学习不是在训练完成后添加的装饰品——它们参与了定义 KNN 模型本身。你不选出距离怎么定义的，就没有模型。

\subsection{高维中局部需要用指数级样本购买}

在 \([0,1]^p\) 的超立方体里均匀撒 \(n\) 个点。要圈出一个边长为 \(r\) 的局部小立方体，并期望里面平均包含 \(k\) 个样本，大致需要

\[
nr^p\approx k,
\qquad
r\approx\left(\frac{k}{n}\right)^{1/p}.
\]

看着好像还好：只要 \(n\) 够大，\(r\) 就能压缩到任意小。问题是 \(p\) 坐在指数上。\(p=2\) 时 \(r\approx\sqrt{k/n}\)；\(p=10\) 时 \(r\approx(k/n)^{0.1}\)；\(p=100\) 时 \(r\approx(k/n)^{0.01}\)。

固定 \(k=30\)，\(n=10^6\)。二维下 \(r\approx0.005\)，是个紧贴查询点的小方块。十维下 \(r\approx(3\times10^{-5})^{0.1}\approx0.35\)——所谓``最近邻居''平均分布在一个边长是总范围三分之一的区域里。这和``局部''已经没什么关系了。一百维下 \(r\approx(3\times10^{-5})^{0.01}\approx0.90\)——几乎就是把整个空间都划成邻域。

要保持同一局部半径，样本量必须随维度指数增长。维数灾难对 KNN 的刀不是藏在复杂的定理里，就摆在这个幂指数关系里。

\subsection{距离集中让远近失去意义}

高维里还有第二个效应。许多分布中，各坐标的距离贡献近似独立地相加。中心极限定理暗示，不同点之间的距离会簇拥在一个平均值附近，相对波动随维度缩小。最远点和最近点到查询点的距离之比趋近于一。

这是比``邻域不再局部''更让 KNN 难堪的局面：所有点都差不多远时，最近邻是纯粹由噪声随机拣出来的几个点。在小扰动下最近邻集合可能剧烈变化。

当然，这不等于任意高维 KNN 必然失败。数据如果实际趴在某个低维流形上，有效维度远小于坐标数，邻域仍可能有意义。稀疏表示或合适的度量也能恢复距离的区分力。但成功来自数据结构对 KNN 假设的成全，不是 KNN 绕过了维数灾难。先从不做降维的全坐标距离开始盲目跑 KNN，高维下得到的结果和随机猜没有多少方法论上的差别。

实际中常有人先对整个数据集做 PCA 降维，再在降维空间里交叉验证 KNN 的超参数。这么做会在降维步骤中把测试集的信息泄漏进训练过程，交叉验证给出的误差估计偏乐观。降维必须嵌入每个交叉验证折内部，对训练集学习投影，再对验证集应用同一投影。

\subsection{KNN 不会告诉你它不知道自己不知道}

分类 KNN 的概率估计就是邻域内的类别频率。样本不平衡时，多数类凭数量优势占据邻域中心位置，少数类的样本永远被淹没。距离加权投票或代价敏感阈值可以部分缓解，但代价是改变了被估计的量的定义。

更致命的是查询点落在训练数据的支撑集之外时。KNN 没有任何机制发出警告。它会忠实地返回某几个``最近''的点，即使最近的那个也远在天边。在输入空间的高维空洞里，算法仍然给你一个预测，让你误以为自己有了答案。

\subsection{查询成本与近似妥协}

一次朴素近邻查询需要扫过全部 \(n\) 个训练样本，计算 \(p\) 维距离——\(O(np)\)。KD-tree、ball-tree 等空间索引在低维下能把查询压到 \(O(\log n)\)，但维度升高后多数树结构退化成线性扫描。近似最近邻方法以可控的召回误差换速度，但引入了一层与模型本身无关的误差源。评价 KNN 系统时应该把搜索近似造成的误差和统计模型的估计误差分开算账，而不是用一个端到端的数字掩盖各自的贡献。

KNN 的最大价值也许不是作为生产系统里的预测器，而是把归纳偏置暴露得无比直白：相近的输入应该有相近的输出。困难同样直白——相近由谁定义，以及你手中的样本量是否足够让局部真正成为局部。后面的核方法、高斯过程乃至深度学习里的注意力，都可以看作这一直觉在不同层面的精细变体：它们都要回答``谁和谁相近''，只是不再用赤裸裸的 Euclidean 距离和原始坐标。
